% Chapter 6 of Singular Learning Theory for the Rest of Us: Newton polygons.
\sltchapter{Newton Polygons}{The toric shortcut}{$y=a^2x_1+b^3x_2$}{ch:newton}
\newtonfoot{0/6,4/0}{12/5}{\tfrac5{12}}{a^4+b^6}
\begin{refsection}[chapters/06-newton.bib]

\begin{chaptersummary}
Can the learning coefficient be read off the exponents of the loss? For the model $y=a^2x_1+b^3x_2$, whose loss is $K=\frac12(a^4+b^6)$, logarithmic coordinates turn the set of nearly true parameters into a polyhedron and its volume exponent into the value of a linear program, $\frac5{12}$, whose dual finds the point $(t,t)=(\frac{12}5,\frac{12}5)$ where the diagonal meets the Newton polygon of $K$. We prove that $\lambda=1/t$ for every sum of even monomials with positive coefficients in two variables, with multiplicity $m=2$ exactly when the diagonal meets a vertex, compute the zeta function on the cones of the polygon's normal fan, and state with sources the theorem for nondegenerate functions of $d$ variables. Exact quadrature for thirty random sums reproduces Newton's $\lambda$ to $3\times10^{-7}$ and his $m$ to $10^{-3}$, while Monte Carlo slopes at practical $\eps$ fall short by about $0.09$ when $m=2$. The polygon depends on the coordinates: the losses of the product model at a generic truth and of a linear network with two hidden units are degenerate, and for them $1/t$ is only an upper bound.
\end{chaptersummary}

\section{The question: reading \texorpdfstring{$\lambda$}{lambda} off the exponents}\label{sec:newton:question}

The learning coefficients computed so far came from changes of coordinates. \Cref{ch:blowups} blew up the origin twice to find $\lambda=\frac34$ for $a^2+b^4$, and three times for the cusp $(a^2-b^3)^2$, where $\lambda=\frac5{12}$; each computation ended with the loss a monomial times a positive function in every chart, and with $\lambda$ read off the monomial's exponents (\cref{thm:blowups:pullback}). Do the exponents of the loss itself, before a change of coordinates, already contain the answer? The loss of a linear network with Gaussian inputs is a polynomial in the weights, and so are the losses of the toy models of this book: their exponents are what one sees first.

The model of this chapter is the \term{mixed cup}\index{mixed cup} $y=a^2x_1+b^3x_2+\varepsilon$, with independent standard normal inputs $x_1,x_2$, unit noise, the uniform prior on $[-1,1]^2$, and the truth $(a,b)=(0,0)$, so that $y$ is pure noise. Its loss is
\[
K(a,b)=\tfrac12\E_x\bigl[(a^2x_1+b^3x_2)^2\bigr]=\tfrac12\bigl(a^4+b^6\bigr),
\]
whose Hessian at the truth is zero: the data see $a$ only through $a^2$ and $b$ only through $b^3$.

\begin{regularanswer}[exponents]
If $K$ has a nondegenerate minimum at $w_0$, then in suitable analytic coordinates centred at $w_0$, $K=\sum_{i=1}^dw_i^2$ (the Morse lemma, the case $r=d$ of \cref{thm:toolkit:splitting}). Every exponent that matters is $2$, and $\lambda=\frac d2$ with $m=1$ (\cref{ch:regular}).
\end{regularanswer}

For the mixed cup the regular answer says nothing: there are no squares, and counting parameters gives $\frac d2=1$. Since $K$ is a sum of a function of $a$ and a function of $b$, the rule for sums of functions of separate variables (\cref{prop:blowups:rules}) gives $\lambda=\frac14+\frac16=\frac5{12}$ and $m=1$: each coordinate contributes the reciprocal of its exponent. That rule stops at sums of separate monomials. The losses $a^6+a^2b^2+b^6$ and $a^4b^2+a^2b^4$, the emblems of Parts~III and~IV, mix the variables.\sidenote{Each opening page of this book carries the Newton diagram of its emblem, and each part page that of its part; this chapter explains the drawing.}

The learning coefficient of a function is what algebraic geometers call its \term{real log canonical threshold}\index{real log canonical threshold}; the two names denote one number. The question of the chapter is: \emph{given the exponents of the monomials of $K$, what are $\lambda$ and $m$, and when is the answer read off the exponents the true one?}\index{learning coefficient!from exponents} The answer was found in the 1970s for the decay of oscillatory integrals,\sidenote{An oscillatory integral $\int e^{i\tau S(w)}\phi(w)\dd w$ decays like a power of $\tau$, possibly times a power of $\log\tau$, as $\tau\to\infty$. For a phase $S\ge0$ it obeys the same Newton-polyhedron estimates as the volume of $\{S<\eps\}$ \parencite[Theorem~1.6]{Greenblatt2010}.} and it is drawn at the foot of the chapter's first page: mark the exponents of $K$, take the convex hull of the quadrants above them, and find where the diagonal enters it. For the mixed cup it enters at $(t,t)=(\frac{12}5,\frac{12}5)$, and $1/t=\frac5{12}$.

\section{By hand: Newton diagrams of the cast}\label{sec:newton:byhand}

Write a polynomial as $K=\sum_{v\in\cE}c_va^{v_1}b^{v_2}$, with $\cE$ the set of exponents $v$ whose coefficients $c_v$ are not zero. The \term{Newton polygon}\index{Newton polygon} $\Gamma(K)$ is the convex hull of the quadrants $v+\R^2_{\ge0}$, $v\in\cE$. Its boundary comes down a vertical ray, runs through finitely many vertices and edges, and leaves along a horizontal ray. The diagonal enters $\Gamma(K)$ at a point $(t,t)$, and $t$ is the \term{Newton distance}\index{Newton distance}. The rule this chapter proves and tests is
\begin{equation}\label{eq:newton:rule}
\lambda=\frac1t,\qquad m=\begin{cases}2&\text{if $(t,t)$ is a vertex of $\Gamma(K)$,}\\1&\text{otherwise.}\end{cases}
\end{equation}
\Cref{fig:newton:diagrams} draws three Newton diagrams, and \cref{tab:newton:cast} applies the rule to the cast of the book. It reproduces every value computed in \crefrange{ch:coefficient}{ch:blowups}, and it fails on three of the last four rows.

\begin{figure}[t]
\begin{fullwidth}
\centering
\setlength{\tabcolsep}{5mm}
\begin{tabular}{ccc}
\newtondiagram{2/4}{4}{\tfrac14}{a^2b^4} &
\newtondiagram{0/6,2/2,6/0}{2}{\tfrac12}{a^6+a^2b^2+b^6} &
\newtondiagram{0/4,2/0}{4/3}{\tfrac12\neq\tfrac1t}{(a-b^2)^2}
\end{tabular}
\end{fullwidth}
\caption{Newton diagrams. Grey dots are the exponents $(i,j)$ of $a^ib^j$, $0\le i,j\le6$; the shaded region is the Newton polygon, its bronze boundary passes through the vertices, the dashed slate line is the diagonal, and the ringed point is $(t,t)$. For $a^2b^4$ the diagonal meets the horizontal ray at $(4,4)$; for $a^6+a^2b^2+b^6$ it meets the vertex $(2,2)$, and $m=2$. The polygon of $(a-b^2)^2=a^2-2ab^2+b^4$ is that of $a^2+b^4$, but its learning coefficient is $\frac12$, not $1/t=\frac34$.}\label{fig:newton:diagrams}
\end{figure}

\begin{table}[t]
\centering\small
\setlength{\tabcolsep}{4pt}
\begin{tabular}{@{}llccc@{}}
\toprule
model & $K$ & vertices of $\Gamma(K)$ & Newton $(\frac1t,m)$ & true $(\lambda,m)$\\
\midrule
regular plane & $\frac12(a^2+b^2)$ & $(0,2),(2,0)$ & $(1,1)$ & $(1,1)$\\
product model & $\frac12a^2b^2$ & $(2,2)$ & $(\frac12,2)$ & $(\frac12,2)$\\
cusp model & $\frac12a^2b^4$ & $(2,4)$ & $(\frac14,1)$ & $(\frac14,1)$\\
quartic model & $\frac12(a^2+b^4)$ & $(0,4),(2,0)$ & $(\frac34,1)$ & $(\frac34,1)$\\
quartic cup & $\frac12(a^4+b^4)$ & $(0,4),(4,0)$ & $(\frac12,1)$ & $(\frac12,1)$\\
mixed cup & $\frac12(a^4+b^6)$ & $(0,6),(4,0)$ & $(\frac5{12},1)$ & $(\frac5{12},1)$\\
emblem of Part III & $a^6+a^2b^2+b^6$ & $(0,6),(2,2),(6,0)$ & $(\frac12,2)$ & $(\frac12,2)$\\
emblem of Part IV & $a^4b^2+a^2b^4$ & $(2,4),(4,2)$ & $(\frac13,1)$ & $(\frac13,1)$\\
\midrule
& $(a-b)^2$ & $(0,2),(2,0)$ & $(1,1)$ & $(\frac12,1)$\\
& $(a-b^2)^2$ & $(0,4),(2,0)$ & $(\frac34,1)$ & $(\frac12,1)$\\
& $(a^2-b^2)^2$ & $(0,4),(4,0)$ & $(\frac12,1)$ & $(\frac12,2)$\\
cusp curve & $(a^2-b^3)^2$ & $(0,6),(4,0)$ & $(\frac5{12},1)$ & $(\frac5{12},1)$\\
\bottomrule
\end{tabular}
\caption{The rule \eqref{eq:newton:rule} on the cast. Upper block: sums of even monomials with positive coefficients, where \cref{thm:newton:plane} proves the rule; the true values come from \cref{ch:coefficient,ch:volume,ch:blowups} and, for the two emblems, from the theorem. Lower block: functions with cancellations; the true values come from \cref{comp:newton:degenerate} and \cref{thm:newton:greenblatt}. Script: \texttt{figures/scripts/06-newton-byhand.py}.}\label{tab:newton:cast}
\end{table}

\begin{computation}{The mixed cup in logarithmic coordinates}\label{comp:newton:emblem}
\step The prior is uniform on $[-1,1]^2$ and $K=\frac12(a^4+b^6)$ is even in each variable, so the volume of nearly true parameters is $V(\eps)=\mathrm{area}\{(a,b)\in(0,1]^2:a^4+b^6<2\eps\}$.
\step Substitute $a=e^{-u_1}$, $b=e^{-u_2}$ with $u\in[0,\infty)^2$, so that $\dd a\,\dd b=e^{-u_1-u_2}\dd u$, $a^4=e^{-4u_1}$ and $b^6=e^{-6u_2}$. Write $L=\log\frac1\eps$.
\step A sum of two positive numbers lies between the larger and twice the larger. So $a^4+b^6<2\eps$ forces $4u_1>L-\log2$ and $6u_2>L-\log2$, and it is implied by $4u_1>L$ and $6u_2>L$. The set is squeezed between two translated quadrants, whose weights $\int e^{-u_1-u_2}\dd u$ are $e^{-L/4-L/6}=\eps^{5/12}$ and $2^{5/12}\eps^{5/12}$: for $\eps\le\frac12$,
\[
\eps^{5/12}\le V(\eps)\le2^{5/12}\eps^{5/12}=1.3348\,\eps^{5/12}.
\]
\step The exponent is the value of a linear program: in the rescaled coordinates $u/L$ the quadrant is $P=\{4u_1\ge1,\ 6u_2\ge1\}$, the weight $e^{-L(u_1+u_2)}$ is largest where $u_1+u_2$ is smallest, and $\min_P(u_1+u_2)=\frac14+\frac16=\frac5{12}$, at the corner $(\frac14,\frac16)$.
\step The dual program, to maximize $\theta_1+\theta_2$ over $\theta\ge0$ subject to $\theta_1(4,0)+\theta_2(0,6)\le(1,1)$, has the optimum $\theta=(\frac14,\frac16)$ with the same value. Dividing $\theta_1(4,0)+\theta_2(0,6)=(1,1)$ by $\theta_1+\theta_2$ exhibits $(\frac{12}5,\frac{12}5)$ as a convex combination of the exponents: the dual finds the point where the diagonal meets the Newton polygon, and $\lambda=1/t$.
\step Exactly, scaling $a$ by $(2\eps)^{1/4}$ and $b$ by $(2\eps)^{1/6}$ gives $V(\eps)=C_6(2\eps)^{5/12}=1.2662\,\eps^{5/12}$ for $\eps\le\frac12$, with $C_6=\int_0^1(1-v^6)^{1/4}\dd v=\Gamma(\frac16)\Gamma(\frac54)/(6\Gamma(\frac{17}{12}))=0.9486$. The partition function factorizes, and $\bar F(n)=\frac5{12}\log n-\log\bigl(2^{5/12}\Gamma(\frac54)\Gamma(\frac76)\bigr)=\frac5{12}\log n-0.1155$ up to exponentially small terms, as \cref{lem:volume:abelian} predicts. At $n=10^4$ the mixed cup pays $3.722$; the regular plane pays $\log n-\log\frac\pi2=8.759$.
\end{computation}

What one had to notice is that \term{logarithmic coordinates}\index{logarithmic coordinates} turn monomials into linear functions, and a sum of monomials with positive coefficients into something within a constant factor of its largest term.\sidenote{This comparison of the sum with its largest term is the only estimate in the proof of \cref{thm:newton:plane}. It fails as soon as a coefficient is negative, and nondegeneracy is the condition that repairs it.} The set of nearly true parameters becomes a polyhedron, its volume an exponential integral over the polyhedron, and the rate of decay of that integral the value of a \term{linear program}\index{linear programming}, whose dual lives in the plane of exponents (\cref{fig:newton:dual}).

\begin{figure}[t]
\centering
\begin{tikzpicture}[x=40mm,y=40mm,font=\footnotesize]
\begin{scope}
  \fill[singular!13] (0.25,0.62) -- (0.25,0.16667) -- (0.62,0.16667) -- (0.62,0.62) -- cycle;
  \draw[muted!80,->] (0,0) -- (0.68,0) node[right] {$u_1$};
  \draw[muted!80,->] (0,0) -- (0,0.68) node[above] {$u_2$};
  \draw[regular,line width=0.6pt,dash pattern=on 2.5pt off 1.5pt] (0.41667,0) -- (0,0.41667);
  \draw[singular,line width=0.7pt] (0.25,0.62) -- (0.25,0.16667) -- (0.62,0.16667);
  \fill[singular] (0.25,0.16667) circle (1pt);
  \node[singular,anchor=west] at (0.26,0.55) {$4u_1\ge1$};
  \node[singular,anchor=south] at (0.5,0.17) {$6u_2\ge1$};
  \node[regular,anchor=north] at (0.41667,-0.01) {$u_1+u_2=\tfrac5{12}$};
  \node[anchor=north] at (0.3,-0.13) {mixed cup: $m=1$};
\end{scope}
\begin{scope}[xshift=60mm]
  \fill[singular!13] (0.16667,0.62) -- (0.16667,0.33333) -- (0.33333,0.16667) -- (0.62,0.16667) -- (0.62,0.62) -- cycle;
  \draw[muted!80,->] (0,0) -- (0.68,0) node[right] {$u_1$};
  \draw[muted!80,->] (0,0) -- (0,0.68) node[above] {$u_2$};
  \draw[regular,line width=0.6pt,dash pattern=on 2.5pt off 1.5pt] (0.5,0) -- (0,0.5);
  \draw[singular,line width=0.7pt] (0.16667,0.62) -- (0.16667,0.33333) -- (0.33333,0.16667) -- (0.62,0.16667);
  \draw[singular,line width=2pt] (0.16667,0.33333) -- (0.33333,0.16667);
  \node[singular,anchor=west] at (0.18,0.55) {$6u_1\ge1$};
  \node[singular,anchor=south] at (0.5,0.17) {$6u_2\ge1$};
  \node[singular,anchor=south west] at (0.28,0.33) {$2u_1+2u_2\ge1$};
  \node[regular,anchor=north] at (0.5,-0.01) {$u_1+u_2=\tfrac12$};
  \node[anchor=north] at (0.3,-0.13) {$a^6+a^2b^2+b^6$: $m=2$};
\end{scope}
\end{tikzpicture}
\caption{The dual picture. In the coordinates $a=e^{-Lu_1}$, $b=e^{-Lu_2}$ the set where every monomial of $K$ is below $\eps=e^{-L}$ is the shaded polyhedron $P$, with one bronze constraint for each vertex of the Newton polygon. The prior weight $e^{-L(u_1+u_2)}$ is largest on the dashed line $u_1+u_2=1/t$, which touches $P$ at a corner for the mixed cup and along the thick edge for $a^6+a^2b^2+b^6$; there the volume collects equal contributions from a range of scales of length proportional to $L$, and $m=2$.}\label{fig:newton:dual}
\end{figure}

\begin{computation}{A vertex on the diagonal: $a^6+a^2b^2+b^6$}\label{comp:newton:vertex}
\step The point $(2,2)$ lies below the segment from $(0,6)$ to $(6,0)$, so all three exponents are vertices, the diagonal enters the polygon at the vertex $(2,2)$, $t=2$, and the rule predicts $\lambda=\frac12$ and $m=2$.
\step Each monomial is below $K$ and $K$ is below three times the largest, so with $R(\delta)=\{(a,b)\in(0,1]^2:a^6<\delta,\ a^2b^2<\delta,\ b^6<\delta\}$ we have $R(\eps/3)\subseteq\{K<\eps\}\subseteq R(\eps)$.
\step For $a<\delta^{1/3}$ the binding constraint on $b$ is $b<\delta^{1/6}$; for $\delta^{1/3}\le a<\delta^{1/6}$ it is $b<\delta^{1/2}/a$. Hence
\[
\mathrm{area}\,R(\delta)=\delta^{1/3}\delta^{1/6}+\int_{\delta^{1/3}}^{\delta^{1/6}}\frac{\delta^{1/2}}a\dd a=\delta^{1/2}\bigl(1+\tfrac16\log\tfrac1\delta\bigr),
\]
and $(\frac\eps3)^{1/2}(1+\frac16\log\frac3\eps)\le V(\eps)\le\eps^{1/2}(1+\frac16\log\frac1\eps)$: $\lambda=\frac12$ and $m=2$, as predicted.
\step The logarithm counts scales. Every scale $a\in(\eps^{1/3},\eps^{1/6})$ contributes equally, because there $a^2b^2$ dominates and the region is a hyperbolic strip; in \cref{fig:newton:dual} this range is the thick edge of $P$, of length $\frac16$ in $u_1$. Quadrature (\cref{sec:newton:measuring}) gives $V(10^{-8})=3.866\times10^{-4}$, between the bounds $2.456\times10^{-4}$ and $4.070\times10^{-4}$, and $V(\eps)/(\eps^{1/2}\log\frac1\eps)$ equal to $0.210$ at $\eps=10^{-8}$ and $0.174$ at $10^{-50}$, approaching $\frac16$. For the product model $a^2b^2$ the constant is $\frac12$: its strip runs over all scales from $\eps^{1/2}$ to $1$.
\end{computation}

The rule looks only at exponents, so it cannot see cancellations between coefficients.

\begin{computation}{When the shortcut fails}\label{comp:newton:degenerate}
\step In $(a-b)^2=a^2-2ab+b^2$ the exponent $(1,1)$ lies on the segment between $(2,0)$ and $(0,2)$, so the polygon is that of $a^2+b^2$: $t=1$, and the rule predicts $\lambda=1$, the value of a regular model with two parameters.
\step The truth is $\lambda=\frac12$: the change $u=a-b$, $w=b$ has Jacobian $1$ and turns $K$ into $u^2$ (invariance under analytic isomorphisms, \cref{prop:blowups:rules}). On the prior square, exactly, $V(\eps)=\sqrt\eps-\frac\eps4$.
\step In logarithmic coordinates step~3 of \cref{comp:newton:emblem} breaks: on the diagonal $a=b>0$ all three monomials have size $a^2$ and $K=0$. The face polynomial of the edge, $(a-b)^2$, vanishes at points with $a\ne0$ and $b\ne0$.
\step The same happens to $(a-b^2)^2$, whose polygon is that of the quartic model while $u=a-b^2$ turns it into $u^2$, so $\lambda=\frac12<\frac34$; and to $(a^2-b^2)^2=(a-b)^2(a+b)^2$, which is $u^2w^2$ in the coordinates $u=a-b$, $w=a+b$, so that $\lambda=\frac12=1/t$ but $m=2$. The cusp curve $(a^2-b^3)^2=a^4-2a^2b^3+b^6$ has the polygon of the mixed cup and an edge polynomial that vanishes on the curve $a^2=b^3$, yet three blow-ups give $\lambda=\frac5{12}=1/t$ and $m=1$ (\cref{ch:blowups}). Cancellation is a danger, not a verdict; \cref{thm:newton:greenblatt} says when it matters.
\end{computation}

\section{The theorem: Newton's polygon in the plane}\label{sec:newton:theorem}

The computations suggest a proof: compare $K$ with its largest monomial, pass to logarithmic coordinates, and solve a linear program. This section carries it out for the functions on which no cancellation can occur, computes their zeta function on the cones of the normal fan, and extracts the half of the rule that survives cancellations.

\subsection{Newton polygons and linear programs}

\begin{definition}[Newton polygon]\label{def:newton:polygon}
Let $K(a,b)=\sum_{v\in\cE}c_va^{v_1}b^{v_2}$ be a polynomial, or a power series convergent near $0$, whose exponent set $\cE=\{v:c_v\ne0\}\subset\N^2$ does not contain $(0,0)$. Its \emph{Newton polygon} is $\Gamma=\conv(\cE)+\R^2_{\ge0}$ and its \emph{Newton distance} is $t=\min\{\tau\ge0:(\tau,\tau)\in\Gamma\}$. For a nonzero $\eta\in\R^2_{\ge0}$, the function $v\mapsto\ip{v}{\eta}$ attains its minimum over $\Gamma$, which is $\mu(\eta)=\min_{v\in\cE}\ip{v}{\eta}$, on a \emph{face} of $\Gamma$ (a vertex, a compact edge, or one of the two boundary rays), and $\eta$ is an inner normal of that face. The \term{diagonal face}\index{diagonal face} is the face that contains $(t,t)$ in its relative interior. For a compact face $F$ the \emph{face polynomial} is $K_F=\sum_{v\in\cE\cap F}c_va^{v_1}b^{v_2}$. A \term{positive even sum}\index{positive even sum} is a polynomial $K$ as above with every $c_v>0$ and every exponent even.
\end{definition}

The minimum over $\Gamma$ equals the minimum over $\cE$ because a point of $\Gamma$ is a convex combination of exponents plus a vector $w\ge0$, and $\ip{w}{\eta}\ge0$; and $t\ge\frac12$, because every point of $\Gamma$ has coordinate sum at least $1$. A positive even sum is nonnegative, vanishes at $0$, and is a function of $\abs a$ and $\abs b$.

Order the vertices of $\Gamma$ as $V^1,\dots,V^r$ with increasing first coordinates. Let $\eta_0=(1,0)$ and $\eta_r=(0,1)$, the inner normals of the two rays, and for $1\le k<r$ let $\eta_k$ be the primitive integer vector perpendicular to the edge $[V^k,V^{k+1}]$, whose coordinates are positive. These are the rays of the \term{normal fan}\index{normal fan}; $\cC_k$ denotes the cone spanned by $\eta_{k-1}$ and $\eta_k$, $\abs\eta_1=\eta_1+\eta_2$ the coordinate sum of $\eta\in\R^2_{\ge0}$, and $\mu_k=\mu(\eta_k)$.

\begin{lemma}[The normal fan and the linear program]\label{lem:newton:lp}
Let $\Gamma$ and $t$ be as in \cref{def:newton:polygon}.
\begin{enumerate}
\item The cones $\cC_1,\dots,\cC_r$ cover $\R^2_{\ge0}$ and meet only along rays; the vertex $V^k$ minimizes $\ip{\cdot}{u}$ over $\Gamma$ exactly when $u\in\cC_k$, so $\mu(u)=\ip{V^k}{u}$ there.
\item (Weak duality.) $\mu(u)\le t\abs u_1$ for every $u\in\R^2_{\ge0}$.
\item (Strong duality.) Equality $\mu(u)=t\abs u_1$ holds exactly on the normal cone $\cN$ of $\Gamma$ at $(t,t)$: $\cN=\cC_k$ if $(t,t)=V^k$, and $\cN$ is the ray through $\eta_k$ if $(t,t)$ lies in the relative interior of the face with inner normal $\eta_k$. On the polyhedron $P=\{u\in\R^2_{\ge0}:\ip{v}{u}\ge1\text{ for every }v\in\cE\}$ the minimum of $\abs u_1$ is $1/t$, attained exactly on $F^*=\{u\in\cN:\abs u_1=1/t\}$: the segment from $\eta_{k-1}/\mu_{k-1}$ to $\eta_k/\mu_k$ in the vertex case, the point $\eta_k/\mu_k$ otherwise.
\item For every ray with $\mu_j>0$, $\abs{\eta_j}_1/\mu_j\ge1/t$, with equality exactly for the rays in $\cN$.
\end{enumerate}
\end{lemma}

\begin{proof}
The plan is to read (i) off the convexity of the boundary and (ii)--(iv) off the fact that $(t,t)$ lies on it. (i) The edge from $V^k$ to $V^{k+1}$ is perpendicular to $\eta_k$ and goes right and down, so $\ip{V^{k+1}-V^k}{u}=g_k(u_1\eta_{k,2}-u_2\eta_{k,1})$ with $g_k>0$; it is positive, zero or negative as the angle of $u$ is smaller than, equal to or larger than that of $\eta_k$. By convexity the angles of $\eta_0,\dots,\eta_r$ increase strictly from $0$ to $\frac\pi2$. Hence the cones cover the quadrant and meet only along rays, and for $u\in\cC_k$ the values $\ip{V^j}{u}$ do not increase up to $j=k$ and do not decrease after it, so $V^k$ minimizes over the vertices, hence over $\Gamma=\conv\{V^j\}+\R^2_{\ge0}$; for $u\notin\cC_k$ a neighbouring vertex does strictly better. (ii) Write $(t,t)=\sum_v\theta_vv+w$ with $\theta$ a probability vector on $\cE$ and $w\ge0$; for $u\ge0$, $t\abs u_1=\ip{(t,t)}{u}\ge\sum_v\theta_v\ip{v}{u}\ge\mu(u)$. (iii) Equality in (ii) says that $(t,t)$ minimizes $\ip{\cdot}{u}$ over $\Gamma$, that is, $u\in\cN$. If $(t,t)=V^k$, then $\cN=\cC_k$ by (i). If $(t,t)$ lies in the relative interior of the face with normal $\eta_k$ and direction $e$, then $(t,t)\pm se\in\Gamma$ for small $s>0$, so every $u\in\cN$ is orthogonal to $e$ and is a multiple of $\eta_k$; and $\eta_k\in\cN$, because $(t,t)$ lies on the face where $\ip{\cdot}{\eta_k}$ is smallest. For $u\in P$, (ii) gives $1\le\mu(u)\le t\abs u_1$, with equality exactly when $u\in\cN$ and $\abs u_1=1/t$, and such $u$ lie in $P$. The rays of $\cN$ meet the line $\abs u_1=1/t$ at $\eta/(t\abs\eta_1)=\eta/\mu(\eta)$. (iv) Apply (ii) and (iii) to $u=\eta_j$.
\end{proof}

Minimizing $\abs u_1$ over $P$ is a linear program whose dual, maximizing $\sum_v\theta_v$ over $\theta\ge0$ with $\sum_v\theta_vv\le(1,1)$, has the value $1/t$; parts (ii) and (iii) are weak and strong duality.

\subsection{The formula}

The proof bounds the volume function from both sides; the next lemma, a two-sided form of the Abelian lemma of \cref{ch:volume}, turns such bounds into the learning coefficient and its multiplicity.

\begin{lemma}[Two-sided volume bounds]\label{lem:newton:bounds}
Let $K\ge0$ be bounded and measurable on $W$, let $\varphi$ be a prior density on $W$, and suppose that for some $\lambda>0$, integer $m\ge1$, constants $0<c\le C$ and $\eps_0\in(0,1)$,
\begin{equation}\label{eq:newton:twosided}
c\,g(\eps)\le V(\eps)\le C\,g(\eps)\quad\text{for }0<\eps<\eps_0,\qquad g(\eps)=\eps^\lambda\bigl(\log\tfrac1\eps\bigr)^{m-1}.
\end{equation}
Then $\bar F(n)=\lambda\log n-(m-1)\log\log n+O(1)$. Moreover $\zeta(z)=\int_WK^z\varphi$ is holomorphic on $\Re z>-\lambda$, and $\zeta(x)(x+\lambda)^m$ stays between two positive constants for real $x\in(-\lambda,-\lambda+\delta)$, for some $\delta>0$. Consequently, if $\zeta$ extends meromorphically to a neighbourhood of $-\lambda$, then $-\lambda$ is a pole of order exactly $m$, and no pole has a larger real part.
\end{lemma}

\begin{proof}
The plan is to write both transforms as integrals of $V$ and compare them with the same integrals of $g$. As in the proof of \cref{thm:toolkit:karamata}, $Z(n)=\int_We^{-nK}\varphi=n\int_0^\infty e^{-ns}V(s)\dd s$. Choose $\eps_1\le\eps_0$ with $\eps_1<e^{-(m-1)/\lambda}$, so that $g$ increases on $(0,\eps_1]$, and let $\nu$ be the measure on $[0,\eps_1]$ with $\nu([0,s))=g(s)$. Part (i) of \cref{thm:toolkit:karamata} gives $J(n)=\int e^{-ns}\nu(\dd s)=n\int_0^{\eps_1}e^{-ns}g(s)\dd s+e^{-n\eps_1}g(\eps_1)\sim\Gamma(\lambda+1)n^{-\lambda}(\log n)^{m-1}$. Splitting $Z(n)$ at $\eps_1$ and using \eqref{eq:newton:twosided} below $\eps_1$ and $V\le1$ above it, $c(J(n)-e^{-n\eps_1}g(\eps_1))\le Z(n)\le CJ(n)+e^{-n\eps_1}$, and taking logarithms gives the expansion of $\bar F(n)=-\log Z(n)$.

For real $x\in(-\lambda,0)$, Tonelli's theorem and $K^x=\int_K^\infty(-x)s^{x-1}\dd s$ give $\zeta(x)=(-x)\int_0^\infty s^{x-1}V(s)\dd s$. The part over $s\ge\eps_0$ is at most $\eps_0^x$. Below $\eps_0$ we use \eqref{eq:newton:twosided}; with $\kappa=x+\lambda$ and $\ell_0=\log\frac1{\eps_0}$, the substitution $s=\eps_0e^{-\rho}$ turns $\int_0^{\eps_0}s^{x-1}g(s)\dd s$ into $\eps_0^\kappa\int_0^\infty e^{-\kappa\rho}(\rho+\ell_0)^{m-1}\dd\rho$, and expanding the binomial,
\[
\int_0^{\eps_0}s^{x-1}g(s)\dd s=\eps_0^\kappa\sum_{j=0}^{m-1}\binom{m-1}j\frac{j!\,\ell_0^{m-1-j}}{\kappa^{j+1}}=\frac{(m-1)!}{\kappa^m}\bigl(1+O(\kappa)\bigr).
\]
Since $-x\to\lambda>0$, $\zeta(x)\kappa^m$ stays between positive constants near $\kappa=0$, and $\zeta(x)<\infty$ for $x>-\lambda$. As $\abs{K^z}=K^{\Re z}$, the integral defining $\zeta$ converges absolutely and locally uniformly on $\Re z>-\lambda$, where $\zeta$ is therefore holomorphic. If $\zeta$ is meromorphic near $-\lambda$ with Laurent series $\sum_{j\ge-p}a_j(z+\lambda)^j$, $a_{-p}\ne0$, then $\zeta(x)(x+\lambda)^p\to a_{-p}$, and the two-sided bound forces $p=m$.
\end{proof}

\begin{theorem}[Newton's polygon formula in the plane]\label{thm:newton:plane}
Let $K$ be a positive even sum with Newton distance $t$, and let $m=2$ if $(t,t)$ is a vertex of its Newton polygon and $m=1$ otherwise. Let the prior $\varphi$ be a density on $W=[-1,1]^2$ that is bounded, and bounded below by a positive constant near $0$. Then:
\begin{enumerate}
\item there are constants $0<c\le C$ and $\eps_0>0$ with $c\,\eps^{1/t}(\log\frac1\eps)^{m-1}\le V(\eps)\le C\,\eps^{1/t}(\log\frac1\eps)^{m-1}$ for $0<\eps<\eps_0$;
\item $\bar F(n)=\frac1t\log n-(m-1)\log\log n+O(1)$ as $n\to\infty$;
\item $\zeta$ is holomorphic on $\Re z>-\frac1t$, and $\zeta(x)(x+\frac1t)^m$ stays between two positive constants as $x\downarrow-\frac1t$; for the uniform prior, $\zeta$ is meromorphic on a half-plane containing $\{\Re z\ge-\frac1t\}$ and $-\frac1t$ is its largest pole, of order $m$ (\cref{prop:newton:zeta}).
\end{enumerate}
The learning coefficient of $K$ on $W$ is therefore $\lambda=1/t$, with multiplicity $m$.
\end{theorem}

The true set of a positive even sum is the origin, an axis, or both axes; the theorem is global on $W$, and the origin determines $\lambda$. Its hypotheses on the prior hold for every continuous prior that is positive at the origin.

\begin{proof}
The plan is that of \cref{comp:newton:emblem}: reduce to the uniform prior on the positive quadrant, compare $K$ with its largest monomial, and compute the volume of the largest monomial's sublevel set in logarithmic coordinates with \cref{lem:newton:lp}. Parts (ii) and (iii) then follow from (i) by \cref{lem:newton:bounds}, except the meromorphy, which is \cref{prop:newton:zeta}. Write $L=\log\frac1\eps$ and, for $H\ge0$, $A_H(\eps)=\mathrm{area}\{(a,b)\in(0,1]^2:H(a,b)<\eps\}$.

\emph{Step 1: reduction.} Let $\varphi\le\Phi$ on $W$ and $\varphi\ge\phi_0>0$ on $[-r_0,r_0]^2$, $0<r_0\le1$. By evenness, $V(\eps)\le4\Phi A_K(\eps)$, and $V(\eps)\ge4\phi_0r_0^2A_{\tilde K}(\eps)$ with $\tilde K(a,b)=K(r_0a,r_0b)$, a positive even sum with the same exponents. So it suffices to show, for every positive even sum, that $A_K(\eps)$ lies between constant multiples of $\eps^{1/t}L^{m-1}$ for small $\eps$.

\emph{Step 2: the largest monomial.} Let $M(a,b)=\max_{v\in\cE}a^{v_1}b^{v_2}$ on $(0,1]^2$, and let $c_+$ and $c_-$ be the largest and the smallest coefficient. Then $c_-M\le K\le\abs\cE c_+M$, and
\begin{equation}\label{eq:newton:sandwich}
A_M\bigl(\eps/(\abs\cE c_+)\bigr)\le A_K(\eps)\le A_M(\eps/c_-).
\end{equation}

\emph{Step 3: logarithmic coordinates.} The map $u\mapsto(e^{-u_1},e^{-u_2})$ is a diffeomorphism of the open quadrant onto $(0,1)^2$ with Jacobian $e^{-\abs u_1}$, and $M(e^{-u_1},e^{-u_2})=e^{-\mu(u)}$. Since $\mu$ is positively homogeneous and $\{\mu=1\}$ is a null set,
\[
A_M(e^{-L})=\int_{\{\mu(u)>L\}}e^{-\abs u_1}\dd u=L^2\int_Pe^{-L\abs u_1}\dd u=L^2\int_{1/t}^\infty e^{-Lh}\psi(h)\dd h.
\]
The last step uses the coordinates $(u_1,h)$, $h=\abs u_1$, of Jacobian $1$, and $\psi(h)$ is the length of the interval $\{u_1:(u_1,h-u_1)\in P\}$, an interval because $P$ is convex, and empty for $h<1/t$ by \cref{lem:newton:lp}(iii). Let $\psi_0=\psi(1/t)$, the length of the projection of $F^*$ to the first axis: $\psi_0>0$ if $(t,t)$ is a vertex and $\psi_0=0$ otherwise.

\emph{Step 4: the shape of $P$ near $F^*$.} We show that for some $B\ge1$ and every $h\ge1/t$,
\begin{equation}\label{eq:newton:psi}
\psi_0+\bigl(h-\tfrac1t\bigr)\le\psi(h)\le B\bigl(\psi_0+h-\tfrac1t\bigr).
\end{equation}
Lower bound: $P$ is closed upwards, because the exponents are $\ge0$. If $u^\pm$ are the endpoints of $F^*$ with $u^-_1\le u^+_1$, the points $u^-+(0,h-\frac1t)$ and $u^++(h-\frac1t,0)$ lie in the convex slice $P\cap\{\abs u_1=h\}$, and the segment between them projects to an interval of length $\psi_0+h-\frac1t$. Upper bound in the vertex case: the slice lies in $\{u\ge0:\abs u_1=h\}$, so $\psi(h)\le h\le\max(1,\frac1{t\psi_0})(\psi_0+h-\frac1t)$. Otherwise $F^*=\{u^*\}$ and $(t,t)$ lies in the relative interior of a face $\Delta$ on which $\ip{\cdot}{u^*}$ attains its minimum $1$. Take $w^\pm=(t,t)\pm\delta e\in\Delta$, with $e$ a direction of $\Delta$ and $\delta>0$; they lie on the line $\{\ip{\cdot}{u^*}=1\}$, which avoids $0$, so they are linearly independent. For $u\in P$, $\ip{w^\pm}{u}\ge\mu(u)\ge1=\ip{w^\pm}{u^*}$, so $\kappa_\pm=\ip{w^\pm}{u-u^*}\ge0$ and $\abs u_1-\frac1t=\frac1{2t}\ip{w^++w^-}{u-u^*}=\frac1{2t}(\kappa_++\kappa_-)$. Expanding $u-u^*$ in the basis dual to $(w^+,w^-)$ gives $\abs{u_1-u_1^*}\le C_\Delta(\kappa_++\kappa_-)=2tC_\Delta(\abs u_1-\frac1t)$, so $\psi(h)\le4tC_\Delta(h-\frac1t)$.

\emph{Step 5: conclusion.} Integrating \eqref{eq:newton:psi}, $\int_{1/t}^\infty e^{-Lh}(\psi_0+h-\frac1t)\dd h=e^{-L/t}(\psi_0L^{-1}+L^{-2})$, so for every $L>0$
\begin{equation}\label{eq:newton:AM}
e^{-L/t}(\psi_0L+1)\le A_M(e^{-L})\le B\,e^{-L/t}(\psi_0L+1).
\end{equation}
By \eqref{eq:newton:sandwich}, $A_K(e^{-L})$ lies between the values of \eqref{eq:newton:AM} at $L+\log(\abs\cE c_+)$ and at $L+\log c_-$; a bounded shift of $L$ changes $e^{-L/t}$ and, for large $L$, $\psi_0L+1$ by bounded factors. So $A_K(\eps)$ lies between constant multiples of $\eps^{1/t}(\psi_0L+1)$, which is of the order of $\eps^{1/t}L^{m-1}$ by Step~3, and Step~1 proves (i).
\end{proof}

For the product model, where $M=K=a^2b^2$ and $\psi(h)=h$, the lower bound in \eqref{eq:newton:AM} is the exact volume $\eps^{1/2}(1+\frac12\log\frac1\eps)$. The multiplicity is one plus the dimension of the face of $P$ on which the linear program attains its optimum.

\subsection{The zeta function on the normal fan}

On the cone $\cC_k$ the largest monomial is $a^{V^k_1}b^{V^k_2}$, and a linear change of variables in logarithmic coordinates that maps the quadrant onto $\cC_k$ turns it into a monomial in new variables, and every other monomial into that monomial times powers of the new variables.

\begin{proposition}[The zeta function on the normal fan]\label{prop:newton:zeta}\index{zeta function!on the normal fan}
Let $K$ be a positive even sum, with the notation of \cref{lem:newton:lp}, and let the prior be uniform on $[-1,1]^2$. Let $D_k=\abs{\det(\eta_{k-1},\eta_k)}$, $A_k(z)=\mu_{k-1}z+\abs{\eta_{k-1}}_1$, $B_k(z)=\mu_kz+\abs{\eta_k}_1$, and
\[
Q_k(\xi)=\sum_{v\in\cE}c_v\,\xi_1^{\ip{v}{\eta_{k-1}}-\mu_{k-1}}\xi_2^{\ip{v}{\eta_k}-\mu_k},
\]
a polynomial with nonnegative integer exponents and constant term $c_{V^k}$. Then for $\Re z>0$
\begin{equation}\label{eq:newton:zetafan}
\zeta(z)=\sum_{k=1}^rD_k\int_0^1\!\!\int_0^1\xi_1^{A_k(z)-1}\xi_2^{B_k(z)-1}Q_k(\xi)^z\dd\xi_1\dd\xi_2,
\end{equation}
and $\zeta$ extends meromorphically to the half-plane $\Re z>-\min_j(\abs{\eta_j}_1+1)/\mu_j$ (minimum over $j$ with $\mu_j>0$), which contains $\{\Re z\ge-\frac1t\}$. There its poles lie among the numbers $-\abs{\eta_j}_1/\mu_j$, and the largest is $-\frac1t$, of order $m$. If $(t,t)=V^k$, then $(z+\frac1t)^2\zeta(z)\to D_kc_{V^k}^{-1/t}/(\mu_{k-1}\mu_k)=\frac{\psi_0}tc_{V^k}^{-1/t}$ as $z\to-\frac1t$; otherwise the pole is simple with a positive residue.
\end{proposition}

\begin{proof}
The plan is to change variables on each cone, to continue each integral by splitting off the values of the integrand on two edges of the unit square, and to locate the largest pole with \cref{lem:newton:lp}. By evenness, $\zeta(z)=\int_{(0,1]^2}K^z\dd a\dd b$. The map $\Psi_k(\xi)=(\xi_1^{\eta_{k-1,1}}\xi_2^{\eta_{k,1}},\xi_1^{\eta_{k-1,2}}\xi_2^{\eta_{k,2}})$, the composite of $\xi\mapsto-\log\xi$, the linear map $s\mapsto s_1\eta_{k-1}+s_2\eta_k$ and $u\mapsto e^{-u}$, maps $(0,1)^2$ diffeomorphically onto the points of $(0,1)^2$ whose logarithmic coordinates lie in the interior of $\cC_k$, with Jacobian $D_k\xi_1^{\abs{\eta_{k-1}}_1-1}\xi_2^{\abs{\eta_k}_1-1}$; by \cref{lem:newton:lp}(i) these images cover $(0,1)^2$ up to a null set. Since $a^{v_1}b^{v_2}\circ\Psi_k=\xi_1^{\ip{v}{\eta_{k-1}}}\xi_2^{\ip{v}{\eta_k}}$, $K\circ\Psi_k=\xi_1^{\mu_{k-1}}\xi_2^{\mu_k}Q_k$; the exponents of $Q_k$ are nonnegative, and both vanish only for $V^k$, the one exponent minimizing both $\ip{\cdot}{\eta_{k-1}}$ and $\ip{\cdot}{\eta_k}$. This proves \eqref{eq:newton:zetafan}, and $c_{V^k}\le Q_k\le\sum_vc_v$ on $[0,1]^2$.

Fix $k$, write $A=A_k$, $B=B_k$, and let $\chi=Q_k^z=\exp(z\log Q_k)$, infinitely differentiable on $[0,1]^2$ and entire in $z$. Split $\chi$ into four parts, the last of which collects what remains:
\[
\chi(\xi)=\chi(0,0)+\bigl[\chi(\xi_1,0)-\chi(0,0)\bigr]+\bigl[\chi(0,\xi_2)-\chi(0,0)\bigr]+\chi_{12}(\xi).
\]
By the mean value theorem the brackets are at most $C\xi_1$ and $C\xi_2$, and $\chi_{12}=\int_0^{\xi_1}\!\int_0^{\xi_2}\partial_1\partial_2\chi$ is at most $C\xi_1\xi_2$, uniformly for $z$ in compact sets. Integrating term by term for $\Re A,\Re B>0$,
\[
\int_{[0,1]^2}\xi_1^{A-1}\xi_2^{B-1}\chi=\frac{\chi(0,0)}{AB}+\frac{I_1}{B}+\frac{I_2}{A}+I_{12},
\]
with $I_1=\int_0^1\xi_1^{A-1}[\chi(\xi_1,0)-\chi(0,0)]\dd\xi_1$, $I_2$ the analogous integral in $\xi_2$, and $I_{12}=\int\xi_1^{A-1}\xi_2^{B-1}\chi_{12}$, all three holomorphic as long as $\Re A>-1$ and $\Re B>-1$. This continues each term to that region, with poles only where $A$ or $B$ vanishes, that is, at numbers $-\abs{\eta_j}_1/\mu_j\le-\frac1t$, those equal to $-\frac1t$ coming from rays in $\cN$ (\cref{lem:newton:lp}(iv)); and the region contains $\{\Re z\ge-\frac1t\}$ because $(\abs{\eta_j}_1+1)/\mu_j>\frac1t$.

If $(t,t)=V^k$, both rays of $\cC_k$ lie in $\cN$, so that $A_k(z)=\mu_{k-1}(z+\frac1t)$ and $B_k(z)=\mu_k(z+\frac1t)$. Then the term $D_k\chi(0,0)/(A_kB_k)$ has a double pole, with the coefficient $D_kc_{V^k}^{-1/t}/(\mu_{k-1}\mu_k)>0$, and every other term has at most a simple pole. Since $\mu_j=t\abs{\eta_j}_1$ for these two rays, the length $\psi_0$ of the projection of the segment $F^*$ is $\abs{\eta_{k-1,1}/\mu_{k-1}-\eta_{k,1}/\mu_k}=tD_k/(\mu_{k-1}\mu_k)$. Otherwise one ray $\eta_k$ lies in $\cN$, shared by at most two cones. In a cone whose first ray is $\eta_k$, $B(-\frac1t)>0$ and the coefficient of $1/A$ is $\chi(0,0)/B+I_2=\int_0^1\xi_2^{B-1}\chi(0,\xi_2)\dd\xi_2$, so the residue there is $\frac{D}{\mu(\eta_k)}\int_0^1\xi_2^{B(-1/t)-1}Q(0,\xi_2)^{-1/t}\dd\xi_2>0$; a cone whose second ray is $\eta_k$ is treated symmetrically. The residues add, and the pole is simple.
\end{proof}

For the mixed cup the rays are $(1,0)$, $(3,2)$ and $(0,1)$, and \eqref{eq:newton:zetafan} reads $\zeta(z)=2^{-z}\bigl[2\int_0^1(1+\xi^4)^z\dd\xi+3\int_0^1(1+\xi^6)^z\dd\xi\bigr]/(12z+5)$: one simple pole, at $-\frac5{12}$, with residue $0.5276$, which is $\frac5{12}$ times the constant $1.2662$ of \cref{comp:newton:emblem}, as the correspondence of \cref{thm:volume:zeta} requires. For $a^6+a^2b^2+b^6$ the double pole has the coefficient $\frac1{12}$, which corresponds to $V(\eps)\sim\frac16\eps^{1/2}\log\frac1\eps$, the constant approached in \cref{comp:newton:vertex}. The computation works for every prior that is the restriction to $W$ of a smooth function with $\varphi(0)>0$, with $\chi$ multiplied by the sum of the prior over the four quadrants; iterating the integration by parts as after \cref{def:toolkit:mellin} continues $\zeta$ to all of $\C$.

The maps $\Psi_k$ are the charts of a \term{toric modification}\index{toric modification}, one monomial map for each cone of the normal fan.\sidenote{Toric varieties are named after the torus $(\C\setminus\{0\})^d$, which they contain as a dense open set; here the real torus $(\R\setminus\{0\})^d$, where no coordinate vanishes, plays its part.} When every cone is \emph{smooth}, $D_k=1$, the charts glue to a manifold on which $K$ is a monomial times a positive function: a resolution of singularities in the sense of \cref{thm:toolkit:hironaka}, found without a blow-up chosen by hand and without the existence theorem of \textcite{Hironaka1964a}.\sidenote{For $a^2+b^4$ the fan has the rays $(1,0)$, $(2,1)$ and $(0,1)$; the cone spanned by $(2,1)$ and $(0,1)$ has $D=2$, and adding the ray $(1,1)$ makes every cone smooth. The three charts of the refined fan are those of the two point blow-ups of \cref{ch:blowups}.} For \eqref{eq:newton:zetafan} smoothness is not needed, because the change of variables happens inside the open quadrant.

\subsection{The half that survives cancellation}

The lower bound in \eqref{eq:newton:AM} used only the exponents, and it survives cancellations, because a power series is never much larger than the largest of its vertex monomials.

\begin{corollary}[The Newton bound]\label{cor:newton:bound}
Let $K\ge0$ be real analytic near $0\in\R^2$, $K(0)=0$, $K\not\equiv0$; let $t$ be the Newton distance of its Taylor series and $m_\Gamma$ the multiplicity that \eqref{eq:newton:rule} assigns to its polygon; and let the prior be bounded below by a positive constant near $0$. Then $V(\eps)\ge c\,\eps^{1/t}(\log\frac1\eps)^{m_\Gamma-1}$ for small $\eps$. Hence $\lambda(0)\le1/t$, and $m(0)\ge m_\Gamma$ if $\lambda(0)=1/t$.
\end{corollary}

\begin{proof}
The plan is to bound $K$ by a multiple of its largest vertex monomial and to apply the lower half of the proof of \cref{thm:newton:plane}. Let $\cV$ be the vertex set of $\Gamma$, $M(w)=\max_{v\in\cV}\abs a^{v_1}\abs b^{v_2}$, $\bar v=\max_{v\in\cV}\abs v_1$, and choose $r\in(0,\frac12]$ such that the Taylor series $\sum_\alpha c_\alpha w^\alpha$ converges absolutely on $[-2r,2r]^2$ and the prior is at least $\phi_0>0$ on $[-r,r]^2$. Each exponent $\alpha$ lies in $\Gamma$, so $\alpha\ge\sum_i\theta_iV^i$ coordinatewise for a probability vector $\theta$ on $\cV$. For $\abs a,\abs b\le2r$ the entries of $w/(2r)$ are at most $1$ in absolute value; dropping the factors of $\alpha-\sum_i\theta_iV^i$ and bounding a weighted geometric mean by its largest factor,
\[
\abs{w^\alpha}=(2r)^{\abs\alpha_1}\Bigl|\Bigl(\frac w{2r}\Bigr)^\alpha\Bigr|\le(2r)^{\abs\alpha_1}\prod_i\Bigl|\Bigl(\frac w{2r}\Bigr)^{V^i}\Bigr|^{\theta_i}\le(2r)^{\abs\alpha_1-\bar v}M(w).
\]
Summing, $K\le C_KM$ on $[-2r,2r]^2$, with $C_K=(2r)^{-\bar v}\sum_\alpha\abs{c_\alpha}(2r)^{\abs\alpha_1}$. Substituting $w=rw'$ and using $M(rw')\le M(w')$, $V(\eps)\ge4\phi_0r^2A_M(\eps/C_K)$, and Steps~3--5 of the proof of \cref{thm:newton:plane}, with only the lower bound in \eqref{eq:newton:psi}, give $A_M(e^{-L})\ge e^{-L/t}(\psi_0L+1)$; they used neither signs nor parities. For analytic $K$ the local volume satisfies $V(\eps)\sim C'\eps^{\lambda(0)}(\log\frac1\eps)^{m(0)-1}$ (\cref{thm:volume:zeta}; the meromorphy of $\zeta$ behind it goes back to \textcite{Atiyah1970} and \textcite{BernsteinGelfand1969}), and comparing rates gives the rest.
\end{proof}

\section{The theorem: nondegenerate and degenerate functions}\label{sec:newton:general}

The proof of \cref{thm:newton:plane} used positivity once, in the lower bound $K\ge c_-M$. What it needs is that cancellations among the monomials of a face never make $K$ much smaller than its largest monomial.

\begin{definition}[Nondegenerate functions]\label{def:newton:nondegenerate}
Let $K\ge0$ be real analytic near $0\in\R^d$ with $K(0)=0$, and let $\Gamma(K)=\conv(\cE)+\R^d_{\ge0}$ be the Newton polyhedron of its Taylor series. The function $K$ is \term{nondegenerate}\index{nondegenerate function} if for every compact face $F$ of $\Gamma(K)$ the face polynomial $K_F=\sum_{v\in\cE\cap F}c_vw^v$ has no zero in the torus $(\R\setminus\{0\})^d$.
\end{definition}

Positive even sums are nondegenerate. Varchenko's definition asks instead that $\nabla K_F$ have no zero on the torus; for $K\ge0$ the two agree, because $K(\tau^{\eta_1}w_1,\dots,\tau^{\eta_d}w_d)=\tau^{\mu(\eta)}(K_F(w)+o(1))$ as $\tau\downarrow0$ for $\eta$ normal to $F$ alone, so $K_F\ge0$, and a zero of a nonnegative smooth function is a critical point. The Newton distance and the diagonal face are defined as in the plane.

\begin{theorem}[Newton polyhedra in $d$ variables]\label{thm:newton:general}
Let $K\ge0$ be real analytic near $0\in\R^d$, $K(0)=0$, $K\not\equiv0$; let $t$ be its Newton distance and $k$ the dimension of its diagonal face; and let the prior be smooth near $0$ with $\varphi(0)>0$.
\begin{enumerate}
\item In every analytic coordinate system centred at $0$, $V(\eps)\ge c\,\eps^{1/t}(\log\frac1\eps)^{d-k-1}$ for small $\eps$; so $\lambda(0)\le1/t$, and $m(0)\ge d-k$ if $\lambda(0)=1/t$ \parencite[Theorem~1.2(a)]{Greenblatt2010}.
\item If $K$ is nondegenerate, $V(\eps)$ lies between constant multiples of $\eps^{1/t}(\log\frac1\eps)^{d-k-1}$, and $\lambda(0)=1/t$, $m(0)=d-k$ \parencite[Theorem~1.2(a,b)]{Greenblatt2010}. For ideals: if $K\asymp\sum_if_i^2$ and $\langle f_1,\dots,f_s\rangle$ is monomial or, more generally, sos-nondegenerate, then $(2\lambda(0),m(0))$ is given by the Newton polyhedron of the ideal \parencite[Theorem~1.3]{Lin2017}.
\end{enumerate}
\end{theorem}

In two variables $d-k$ is $2$ for a vertex and $1$ for an edge, which is \eqref{eq:newton:rule}, and (i) is \cref{cor:newton:bound}. Greenblatt's estimates concern $\int_{\{\abs S<\eps\}}\phi$ with a smooth cutoff $\phi\ge0$, $\phi(0)>0$, the local volume function of $S=K$. Varchenko's oscillatory-integral theorem also assumes $t>1$, but the sublevel estimates never need it, and the oscillatory ones hold without it when the phase has a sign \parencite[Theorem~1.6(a)]{Greenblatt2010}: for losses it plays no role.

\begin{proofsketch}
The plan is \cref{prop:newton:zeta} in every dimension, with nondegeneracy in place of positivity. (1) Refine the normal fan of $\Gamma(K)$ to a fan of simplicial cones whose generators are bases of $\Z^d$. (2) A $d$-dimensional cone with generators $\eta^1,\dots,\eta^d$ gives the chart $w_j=\pm\prod_i\xi_i^{\eta^i_j}$; the charts glue to an analytic manifold with a proper map onto a neighbourhood of $0$ that is an isomorphism off the coordinate hyperplanes. (3) In each chart $K=\xi_1^{\mu(\eta^1)}\cdots\xi_d^{\mu(\eta^d)}Q(\xi)$, and the Jacobian is $\prod_i\abs{\xi_i}^{\abs{\eta^i}_1-1}$. On $\{\xi_i=0:i\in I\}$, $Q$ is, up to a monomial, the face polynomial of the face dual to the cone spanned by $\{\eta^i:i\in I\}$, and the points over the origin are those for which that face is compact; nondegeneracy keeps $Q$ away from $0$ there \parencite[Theorem~4.8]{Lin2017}, so every chart is in normal-crossing form. (4) \Cref{thm:blowups:pullback} gives $\lambda$ as the minimum of $\abs\eta_1/\mu(\eta)$ over generators and $m$ as the largest number of minimizing generators in one cone. (5) The linear program of \cref{lem:newton:lp} in $d$ variables (\cref{ex:newton:lp3}) shows that the minimum is $1/t$, attained exactly in the normal cone of the diagonal face, a cone of dimension $d-k$; its refinement supplies $d-k$ minimizing generators in one cone, and no cone has more, because they are linearly independent. Omitted: smooth refinements, properness, and $Q$ over noncompact faces; see \textcite[Section~4]{Lin2017}, who credits the desingularization to \textcite{Varchenko1976} and \textcite[Section~8.3]{ArnoldEtAl1988}, and \textcite[Section~3]{Watanabe2024} for the statistical setting.
\end{proofsketch}

Degenerate functions are where the coordinates matter. In two variables Greenblatt characterized exactly when the Newton distance is still right; the order of a zero of a polynomial is the smallest order of a partial derivative that does not vanish there.

\begin{theorem}[Degenerate functions of two variables; Greenblatt]\label{thm:newton:greenblatt}
Let $K\ge0$ be real analytic near $0\in\R^2$, $K(0)=0$, $K\not\equiv0$, with Newton distance $t$ and diagonal face $\Delta$, and the prior as in \cref{thm:newton:general}.
\begin{enumerate}
\item $\lambda(0)=1/t$, unless $\Delta$ is a compact edge whose face polynomial $K_\Delta$ has a zero in $(\R\setminus\{0\})^2$ of order larger than $t$; then $\lambda(0)<1/t$.
\item If $\lambda(0)=1/t$, then $m(0)=2-\dim\Delta$, unless $K_\Delta$ has a zero in $(\R\setminus\{0\})^2$ of order exactly $t$; then $m(0)=2$.
\end{enumerate}
\end{theorem}

This is Theorem~1.5 of \textcite{Greenblatt2010}, whose multiplicity is our $m-1$, and it settles the lower block of \cref{tab:newton:cast}. The edge polynomials $(a-b)^2$, $(a-b^2)^2$, and $(a+b)^2$ of $(a+b)^2+b^4$ (whose $\lambda$ is $\frac34$, by the substitution $u=a+b$) vanish to order $2$ on curves in the torus, and $2$ exceeds $t=1$, $\frac43$ and $1$. The edge polynomial of $(a^2-b^2)^2$ vanishes to order $2=t$, so $\lambda=\frac12$ and $m=2$. That of the cusp curve vanishes to order $2<t=\frac{12}5$, so the rule is right although $K$ is degenerate. Equivalently, $f=a^2-b^3$ has only zeros of order $1$ on its edge, and $1$ does not exceed its Newton distance $\frac65$; by Greenblatt's Theorem~1.2, $\mathrm{area}\{\abs f<\delta\}\asymp\delta^{5/6}$, and $\{K<\eps\}=\{\abs f<\eps^{1/2}\}$ gives $\lambda=\frac5{12}$.

A degenerate polygon can often be repaired by a change of coordinates, as in \cref{comp:newton:degenerate}. Varchenko proved that every real analytic function of two variables has \term{adapted coordinates}\index{adapted coordinates}, in which the reciprocal of the Newton distance gives the exact decay rate of the associated oscillatory integral (see the account of \textcite[Section~1]{Greenblatt2010}, and \textcite{IkromovMuller2011} for smooth functions). The other repair keeps the coordinates and changes the generators. \textcite[Proposition~3.6]{Lin2017} proves that \emph{if two sets $\{f_1,\dots,f_r\}$ and $\{g_1,\dots,g_s\}$ of functions generate the same ideal of the ring of functions real analytic on a compact semianalytic set $\Omega$, then $\sum_if_i^2$ and $\sum_jg_j^2$ have the same learning coefficient and multiplicity on $\Omega$} (his real log canonical threshold of an ideal is $2\lambda$ in our units). The proof is one line: each $g_j=\sum_ih_{ji}f_i$ with $h_{ji}$ bounded on $\Omega$, so $\sum_jg_j^2\le C\sum_if_i^2$ by the Cauchy--Schwarz inequality, symmetrically, and comparable functions have the same $\lambda$ and $m$ (\cref{prop:blowups:rules}). It makes $\lambda$ a property of the ideal, not of the generators one writes down, and Lin builds Newton polyhedra of ideals on it. \Cref{thm:newton:plane} is the two-variable case of his formula for monomial ideals, since a positive even sum is comparable with the sum of the squares of the monomials $a^{v_1/2}b^{v_2/2}$, $v\in\cE$.

\section{Measuring it}\label{sec:newton:measuring}

Two scripts test the rule on functions whose volume can be computed independently, with the uniform prior on $[-1,1]^2$, so that $V(\eps)$ is the area of $\{K<\eps\}$ in $(0,1]^2$.

\texttt{06-newton-random.py} draws thirty random positive even sums, with two to five distinct exponents from $\{0,2,4,6\}^2\setminus\{(0,0)\}$ and coefficients log-uniform on $[0.2,5]$ (generator seed $6$); eight of them have a vertex on the diagonal. The script computes $t$ and $m$ in rational arithmetic and measures $V(\eps)$ in two ways. The first is quadrature with no input from the theory: since $K$ increases in $b$, $V(\eps)=\int_0^1\beta_\eps(a)\dd a$ with $\beta_\eps(a)=\sup\{b\le1:K(a,b)<\eps\}$; in the coordinate $x=-\log a$ the root is found by Newton's method and the integral by adaptive Gauss--Kronrod quadrature with relative tolerance $10^{-12}$. Logarithmic coordinates reach $\eps=e^{-4000}$ without underflow,\sidenote{That is $\eps\approx10^{-1737}$. No experiment sees such an $\eps$, but the theorem is a statement about $\eps\to0$, and the quadrature can follow it there.} and the routine reproduces the exact volumes of $a^2b^2$, $a^4+b^6$ and $a^2+b^4$ to $10^{-14}$. The exponent and the multiplicity are fitted by least squares of $\log V$ on $-L$, $\log L$, $1$ and $1/L$ at $L=\log\frac1\eps=250,500,\dots,4000$. The second is what a practitioner would do: $2\times10^6$ Monte Carlo points (seed $100+k$ for the $k$-th function) and the slope of $\log V$ between $\eps=10^{-2}$ and $10^{-4}$, with an exact $95\%$ interval from the Clopper--Pearson interval for the fraction of points below $10^{-2}$ that are also below $10^{-4}$. The functions with cancellations have exact volume formulas, evaluated with $40$-digit arithmetic and fitted in the same way.

\begin{figure}[t]
\centering
\begin{tikzpicture}
\begin{axis}[splot,height=0.5\linewidth,xmin=0,xmax=1.08,ymin=0,ymax=1.1,xlabel={Newton's prediction $1/t$},ylabel={measured exponent},
  legend pos=north west,legend style={font=\scriptsize}]
\addplot[regularpred,domain=0:1.08] {1};
\addplot[theory,domain=0:1.08] {x};
\addplot[measured2,ebars] table[x=newton,y=naive,y error plus=nep,y error minus=nem] {figures/data/06-newton-random.dat};
\addplot[measured] table[x=newton,y=quad] {figures/data/06-newton-random.dat};
\addplot[measured3,mark size=2.2pt] table[x=newton,y=quad] {figures/data/06-newton-degenerate.dat};
\legend{regular: $\frac d2=1$,singular: $\lambda=1/t$,Monte Carlo slope at $10^{-4}\le\eps\le10^{-2}$,quadrature fit at $\eps\le e^{-250}$,functions with cancellations}
\end{axis}
\end{tikzpicture}
\caption{Newton's prediction against measured volume exponents. Dots: thirty random positive even sums, exponent fitted to exact quadrature of $V(\eps)$ at $\log\frac1\eps=250$ to $4000$; they lie on the bronze line to within $3\times10^{-7}$. Circles: the same functions, slope of $\log V$ between $\eps=10^{-2}$ and $10^{-4}$ from $2\times10^6$ Monte Carlo points each (seeds $100$ to $129$), with exact $95\%$ intervals, mostly smaller than the marks; the eight with $m=2$ lie about $0.09$ below the line. Squares: functions with cancellations, exponents fitted to exact volumes: from the left the cusp curve and $(a^2-b^2)^2$ on the line, $(a-b^2)^2$ at $(\frac34,\frac12)$, $(a+b)^2+b^4$ at $(1,\frac34)$, and at $(1,\frac12)$ the functions $(a-b)^2$, $\frac12(a_1b_1+a_2b_2)^2$ and the product model at the truth $ab=\frac14$ in local coordinates. Script: \texttt{figures/scripts/06-newton-random.py}.}\label{fig:newton:random}
\end{figure}

The quadrature fits (\cref{fig:newton:random}) agree with Newton's prediction to within $2.5\times10^{-7}$ in $\lambda$, and the fitted multiplicities are within $0.001$ of Newton's, eight of them $2$; the quadrature does not know the polygon, so this agreement was not arranged. The Monte Carlo slopes deviate from $1/t$ by $-0.013$ on average for the twenty-two functions with $m=1$ (range $-0.063$ to $+0.051$) and by $-0.092$ for the eight with $m=2$ (range $-0.110$ to $-0.061$); only eleven of the thirty $95\%$ intervals contain $1/t$. The intervals are correct about the Monte Carlo error and silent about the bias of a slope taken before the asymptotic regime begins.

\begin{figure}[t]
\centering
\begin{tikzpicture}
\begin{axis}[splot,height=0.5\linewidth,xmode=log,xmin=1,xmax=40,ymin=0,ymax=1.08,xlabel={$\log_{10}(1/\eps)$},
  ylabel={local slope $\mathrm{d}\log V/\mathrm{d}\log\eps$},xtick={1,2,5,10,20,40},xticklabels={1,2,5,10,20,40}]
\addplot[regularpred,domain=1:40] {1};
\addplot[theory] table[x=dec,y=p1] {figures/data/06-newton-slopes-curves.dat};
\addplot[theory,restrict x to domain=1.0986:3.7] table[x=dec,y=p2] {figures/data/06-newton-slopes-curves.dat};
\addplot[theory] table[x=dec,y=p3] {figures/data/06-newton-slopes-curves.dat};
\addplot[theory2] table[x=dec,y=s2] {figures/data/06-newton-slopes-curves.dat};
\addplot[theory2] table[x=dec,y=s3] {figures/data/06-newton-slopes-curves.dat};
\addplot[measured,ebars] table[x=dec,y=m1,y error plus=e1p,y error minus=e1m] {figures/data/06-newton-slopes-mc.dat};
\addplot[measured2,ebars] table[x=dec,y=m2,y error plus=e2p,y error minus=e2m] {figures/data/06-newton-slopes-mc.dat};
\addplot[measured3,ebars] table[x=dec,y=m3,y error plus=e3p,y error minus=e3m] {figures/data/06-newton-slopes-mc.dat};
\node[font=\scriptsize,text=regular,anchor=north west] at (axis cs:1.05,0.98) {regular: $\frac d2=1$};
\node[font=\scriptsize,text=singulardeep,anchor=south east] at (axis cs:39,0.5) {$a^6+a^2b^2+b^6$: $\frac12-1/\log\frac1\eps$};
\node[font=\scriptsize,text=singulardeep,anchor=north east] at (axis cs:39,0.405) {$\frac12(a^4+b^6)$: $\frac5{12}$};
\node[font=\scriptsize,text=singulardeep,anchor=north east] at (axis cs:39,0.19) {$a^4b^6+a^6b^4$: $\frac15$};
\end{axis}
\end{tikzpicture}
\caption{What a slope measures: local slopes of $\log V$ against $\log\eps$ for three positive even sums. Solid bronze: the asymptotics of \cref{thm:newton:plane}, $\lambda$ for $m=1$ and $\lambda-1/\log\frac1\eps$ for $m=2$ (drawn from $\eps=10^{-3}$). Dotted bronze: exact local slopes by quadrature (central differences in $\log\eps$); for the mixed cup they coincide with $\frac5{12}$. Marks (dots for $\frac12(a^4+b^6)$, circles for $a^6+a^2b^2+b^6$, squares for $a^4b^6+a^6b^4$): Monte Carlo slopes over the decades from $[10^{-2},10^{-1}]$ to $[10^{-6},10^{-5}]$, $4\times10^6$ uniform points per function (seeds $61$ to $63$), with exact $95\%$ intervals. Script: \texttt{figures/scripts/06-newton-slopes.py}.}\label{fig:newton:slopes}
\end{figure}

\texttt{06-newton-slopes.py} shows where the bias comes from (\cref{fig:newton:slopes}). For the mixed cup the volume is an exact power; the Monte Carlo slopes over five decades are $0.417$, $0.418$, $0.416$, $0.413$ and $0.417$, and four of the five intervals contain $\frac5{12}$. For $a^6+a^2b^2+b^6$ the factor $\log\frac1\eps$ lowers every slope by about $1/\log\frac1\eps$: the exact local slope is $0.386$ at $\eps=10^{-1.5}$, $0.443$ at $10^{-5.5}$ and $0.490$ at $10^{-40}$. For $a^4b^6+a^6b^4$, with $\lambda=\frac15$ and $m=1$, the second pole of $\zeta$, at $-\frac14$, is close, the correction to $V$ is a factor $1-O(\eps^{1/20})$, and the exact slope is $0.141$ at $10^{-3.5}$ and $0.196$ only at $10^{-20}$. Both biases are predicted by the theory, by the multiplicity and by the cones of \cref{prop:newton:zeta}, and both are invisible to a confidence interval; they also lower the Monte Carlo slopes of $(a^2-b^2)^2$, to $0.414$, and of the cusp curve, to $0.382$.

\section{The ledger}\label{sec:newton:ledger}

For sums of even monomials with positive coefficients in two variables, the learning coefficient is read off the exponents, with a complete proof: $\lambda=1/t$, and $m=2$ exactly when the diagonal enters the Newton polygon at a vertex. In every dimension $1/t$ bounds $\lambda$ from above, with equality for nondegenerate functions (stated and sketched).

\ledgercard{The learning coefficient read off the exponents of $K$.}{Every exponent that matters is $2$: $\lambda=\frac d2$, $m=1$.}{$\lambda=1/t$, $m=2$ if the diagonal meets a vertex, else $1$ (\cref{thm:newton:plane}); $\lambda\le1/t$ in every coordinate system (\cref{cor:newton:bound,thm:newton:general}).}{$y=a^2x_1+b^3x_2$, $K=\frac12(a^4+b^6)$: $t=\frac{12}5$, $\lambda=\frac5{12}$, $m=1$, against $\frac d2=1$.}{Thirty random sums: quadrature exponents within $3\times10^{-7}$ of $1/t$, multiplicities within $10^{-3}$; Monte Carlo slopes at $\eps\ge10^{-4}$ low by $0.09$ when $m=2$.}{Nondegeneracy in the chosen coordinates (positive coefficients in the proof); analyticity at the point; a prior positive there.}

\section{Where the theory stops}\label{sec:newton:stops}

\paragraph{What was assumed.} The complete proof assumes positive coefficients and even exponents, a form of nondegeneracy that can be checked by looking. The general theorem assumes nondegeneracy, a property of $K$ in a coordinate system, not of the model: $(a-b)^2$ is degenerate in the coordinates $(a,b)$ and nondegenerate in $(a-b,b)$. It assumes analyticity at the point, since the polygon is drawn from a Taylor series; the loss of a ReLU network is only piecewise analytic, and on a kink there is no Taylor series. It is local: the learning coefficient of a model is the minimum of the local values over its true set (\cref{prop:blowups:rules}), so in general a polygon is needed at every true point. And the prior must be positive at the point: a prior vanishing like $\abs a^{h_1}\abs b^{h_2}$ replaces the diagonal by the ray through $(h_1+1,h_2+1)$ (\cref{ex:newton:prior}).

\paragraph{Where it breaks.} The losses of neural networks are degenerate in their natural coordinates, for a structural reason: a nondegenerate $K\ge0$ vanishes near $0$ only on coordinate hyperplanes (\cref{ex:newton:zeroset} for two variables), and the true sets of networks are not of that form. At a true point $(a_0,b_0)$ of the product model with generic truth $ab=c_0\ne0$, in the coordinates $\alpha=a-a_0$, $\beta=b-b_0$, the loss is $\frac12(b_0\alpha+a_0\beta+\alpha\beta)^2$, whose polygon has the vertices $(2,0)$ and $(0,2)$: the rule predicts the regular value $1$, while $\lambda=\frac12$; the edge polynomial vanishes to order $2>t=1$, as \cref{thm:newton:greenblatt} requires. The linear network $y=(a_1b_1+a_2b_2)x$, with two hidden units, has $K=\frac12(a_1b_1+a_2b_2)^2$; in four variables the diagonal meets the Newton polyhedron at the exponent $(1,1,1,1)$ of the cross term, inside the edge from $(2,2,0,0)$ to $(0,0,2,2)$, and the rule predicts $\lambda=1$ with $m=3$. The truth is $\lambda=\frac12$, $m=1$: $K$ depends on $c=a_1b_1+a_2b_2$, whose density under the uniform prior is continuous at $0$ with the value $\frac14\int_{-1}^1\log^2\abs x\dd x=1$, so \cref{thm:coefficient:free} applies. With one hidden unit the network is the product model and the rule is exact; the second unit makes the true set a quadric cone.\sidenote{The Aoyagi--Watanabe formula of \cref{ch:catalogue} with one input, one output, two hidden units and a truth of rank $0$ gives the same $\lambda=\frac12$ and $m=1$.} Reduced-rank regression behaves the same way: at the zero truth with $M=N=H=2$ the polyhedron of $\frac12\norm{BA}_F^2$ predicts $\lambda=2$, while the formula of \textcite{AoyagiWatanabe2005} gives $\frac32$ (\cref{ex:newton:research}). Even where the rule is exact, slopes at practical $\eps$ are biased (\cref{fig:newton:slopes}), and the same correction, $-(m-1)\log\log n$, affects the estimators of \cref{ch:temperature,ch:llc}.

\paragraph{Open problems.} In two variables adapted coordinates always exist, by Varchenko's theorem \parencite[see][Section~1]{Greenblatt2010}, and \textcite{IkromovMuller2011} extend them to smooth functions. In three or more variables, Greenblatt asks in which situations a coordinate change exists after which the Newton polyhedron gives sharp estimates, and expects the answer to need higher-dimensional resolution of singularities \parencite[Section~1]{Greenblatt2010}. For singular learning theory the question is concrete: for which architectures and truths is there a short, explicit sequence of blow-ups after which the loss is nondegenerate in every chart, so that the rest is a linear program? We know of no general answer, even for two-layer networks with analytic activations. A second question is computational: deciding nondegeneracy means deciding, face by face, whether a polynomial system has a real solution with no zero coordinate.

\begin{exercises}
\item \consolidation Show that adding to a positive even sum a monomial with exponent in its Newton polygon changes neither $\lambda$ nor $m$. For $a^6+b^6$ and $a^6+a^2b^4+b^6$ check with \cref{prop:newton:zeta} that $V(\eps)\sim C\eps^{1/3}$ with $C=0.9638$ and $0.9133$.\label{ex:newton:inside}
\item \consolidation Let $K\ge0$ be analytic near $0\in\R^d$ with $K(0)=0$. Show that every exponent of its Taylor series has coordinate sum at least $2$, so that $t\ge\frac2d$, and conclude from \cref{thm:newton:general}(i) that $\lambda(0)\le\frac d2$. When does the rule \eqref{eq:newton:rule} give $\frac d2$?
\item \consolidation Derive the regular answer from \cref{thm:newton:general}: if $K=\sum_{i=1}^dw_i^2+O(\abs w^3)$, its Newton polyhedron has the vertices $2e_i$, the Newton distance is $\frac2d$, the diagonal face is the simplex spanned by the vertices, and $K$ is nondegenerate.
\item \consolidation\hard Prove \cref{lem:newton:lp}(ii)--(iii) and the bounds of \cref{thm:newton:plane}(i) for positive even sums of $d$ variables, with $m=d-k$ for a diagonal face of dimension $k$. \hint{the face of $P$ where $\abs u_1$ is minimal has dimension $d-k-1$; slice $P$ by the hyperplanes $\abs u_1=h$ and bound the slices' volumes by constant multiples of $(h-\frac1t)^k$, as in \eqref{eq:newton:psi}.}\label{ex:newton:lp3}
\item \byhand Apply the rule to the cusp models $a^2b^4$ and $a^2b^6$ and to the chain $a^2b^2c^2$, say which face the diagonal meets, and recover the values $(\frac14,1)$, $(\frac16,1)$ and $(\frac12,3)$ of \cref{ch:coefficient}.
\item \byhand In three variables find $1/t$ and the dimension of the diagonal face for $a^2b^2+b^2c^2+c^2a^2$, $a^2b^2+c^2$ and $a^2+b^2+c^2$. \hint{the answers are $(\frac34,1)$, $(1,2)$ and $(\frac32,1)$; check the second with the sum rule of \cref{prop:blowups:rules}.}
\item \byhand Show from \eqref{eq:newton:zetafan} that under the uniform prior the zeta function of $a^2+b^4$ is $\bigl[\int_0^1(1+\xi^2)^z\dd\xi+2\int_0^1(1+\xi^4)^z\dd\xi\bigr]/(4z+3)$, and check numerically that $\frac13$ of the bracket at $z=-\frac34$ is the constant $0.8740$ of the exact volume $V(\eps)=0.8740\,\eps^{3/4}$.
\item \byhand For which real $c$ is $K_c=a^4+ca^2b^2+b^4$ nonnegative, and for which nondegenerate? Show that $V(\eps)=C_c\eps^{1/2}$ for small $\eps$, with $C_c=\frac18\int_0^{2\pi}k_c(\theta)^{-1/2}\dd\theta$ and $k_c(\theta)=K_c(\cos\theta,\sin\theta)$, as long as $c>-2$, and explain why $m$ jumps to $2$ at $c=-2$. \hint{$C_1=0.8429$, $C_0=0.9270$, $C_{-1}=1.0783$.}
\item \variation Let the prior be proportional to $\abs a^{h_1}\abs b^{h_2}$ near $0$. Redo the proof of \cref{thm:newton:plane} with the weight $e^{-\ip{h+1}{u}}$, and show that $\lambda=1/t_h$, where $t_h(h_1+1,h_2+1)$ is the point at which the ray through $(h_1+1,h_2+1)$ enters the Newton polygon. For the product model with $h=(1,0)$ show that $\lambda=\frac12$ and $m=1$.\label{ex:newton:prior}
\item \variation Show that a nondegenerate $K\ge0$ of two variables vanishes near $0$ only on the coordinate axes. \hint{in a chart, the face polynomials are the restrictions of $Q$ to the edges $\xi_1=0$ and $\xi_2=0$.} Deduce that the product model at a generic truth is degenerate in the coordinates $(a-a_0,b-b_0)$, and find coordinates in which it is not.\label{ex:newton:zeroset}
\item \measurement Sample $u\in\R^2_{\ge0}$ with density proportional to $e^{-\abs u_1}$ on $\{\abs u_1\ge h_0\}$, $h_0=(L+\log c_-)/t$, and show with \cref{lem:newton:lp}(ii) that $(h_0+1)e^{-h_0}$ times the fraction of samples with $K<\eps$ is unbiased for $V(\eps)$. Use it to estimate the exponents of the thirty functions of \texttt{06-newton-random.py} at $\eps=e^{-100}$ and $e^{-200}$. Which half of \cref{thm:newton:plane} does the sampler take for granted?
\item \measurement Modify \texttt{06-newton-slopes.py} to correct the Monte Carlo slopes of $a^6+a^2b^2+b^6$ for the factor $\log\frac1\eps$, using $m=2$. How much bias remains at $\eps=10^{-3}$, and where does it come from? \hint{$V(\eps)=\eps^{1/2}(c_1\log\frac1\eps+c_0)(1+o(1))$; estimate $c_0/c_1$ from the quadrature.}
\item \research Show that the Newton distance of the reduced-rank regression loss $\frac12\norm{BA}_F^2$ at the zero truth, $B\in\R^{N\times H}$, $A\in\R^{H\times M}$, is $t=2/(H\min(M,N))$. For $M=N=H=2$ compare $1/t=2$ with the value $\frac32$ of \cref{thm:catalogue:aoyagi}; for $H=1$ explain why the rule is exact (\cref{thm:catalogue:rankone}). Then look for a blow-up after which the loss for $M=N=H=2$ is nondegenerate in every chart; we know of no general procedure of this kind for network losses.\label{ex:newton:research}
\end{exercises}

\notesandrefs

Newton polyhedra entered singularity theory in the 1970s: \textcite{Kouchnirenko1976} computed Milnor numbers from them, and \textcite{Varchenko1976} the decay of oscillatory integrals for phases nondegenerate with respect to their Newton polyhedra; the second volume of \textcite{ArnoldEtAl1988} contains Varchenko's theorems (Lin cites its Section~8.3). \textcite{Greenblatt2010} proved the sublevel-set estimates of \cref{thm:newton:general,thm:newton:greenblatt}, starting, as we did, from a decomposition of a neighbourhood of the origin according to which monomials dominate.\sidenote{Greenblatt cuts the neighbourhood into slivers on which the vertex monomials of one face dominate all others (his Section~2); in logarithmic coordinates the slivers become neighbourhoods of the cones of the normal fan.} In complex geometry the log canonical threshold of a monomial ideal is the reciprocal of the Newton distance \parencite{Howald2001}.

In singular learning theory, Newton diagrams were used by \textcite{YamazakiWatanabe2004} for mixtures of binomial distributions and by \textcite{YamazakiEtAl2010} for the Bayesian generalization error, and Watanabe's review notes that for nondegenerate Newton diagrams the resolution is found by a toric modification \parencite[Section~3]{Watanabe2024}. \textcite{Lin2011,Lin2017} developed the theory for statistics: real log canonical thresholds of ideals, the $\tau$-distance for monomial priors, sos-nondegeneracy, and toric resolution. \textcite[Chapter~5]{DrtonEtAl2009} and \textcite[Chapter~17]{Sullivant2018} treat Bayesian integrals of algebraic models in textbook form, and \textcite{Watanabe2009} develops resolution of singularities for learning theory at book length.

The proof of \cref{thm:newton:plane} by logarithmic coordinates and linear programming, \cref{lem:newton:bounds}, and the closed form \eqref{eq:newton:zetafan} are the book's own presentation; the result is the two-variable case of Lin's theorem for monomial ideals \parencite[Theorem~1.3 and Proposition~4.12]{Lin2017}, which he proves by toric resolution. The example $(a+b)^2+b^4$ is Lin's Remark~4.13. Theorem numbers of \textcite{Greenblatt2010} and \textcite{Lin2017} refer to the arXiv versions (v3 of both), which we read.

\paragraph{Translating the literature.} \Textcite{Greenblatt2010} writes $S$ for our $K$, $n$ for the number of variables $d$, $d$ for the Newton distance $t$, $N(S)$ for the Newton polyhedron, $C(S)$ for the diagonal face, $k$ for its dimension, and $I_{\abs S,\phi}(\eps)=\int_{\{\abs S<\eps\}}\phi$ for the local volume function $V(\eps)$; his growth index is our $\lambda$ and his multiplicity our $m-1$. \Textcite{Lin2017} writes $\omega$ for $w$, $\mathcal P(f)$ for the Newton polyhedron, $l$ for $t$ and $\theta$ for $m$; his $\mathrm{RLCT}_\Omega(f;\varphi)$ of a function $f\ge0$ is our $(\lambda,m)$, but for an ideal $\langle f_1,\dots,f_r\rangle$ it is $(2\lambda,m)$ for $K=\sum f_i^2$, because his zeta function is $\int(\sum f_i^2)^{-z/2}\abs\varphi$; his amplitude $\omega^\tau$ is our prior factor $\abs w^\tau$, and his $\tau$-distance is the $t_h$ of \cref{ex:newton:prior} with $h=\tau$. Watanabe writes ``Newton diagram'' for the Newton polygon and ``toric modification'' for the charts $\Psi_k$.

\printchapterbib
\end{refsection}
